\chapter{Cinemática do Contínuo}
\section{Mapeamento do Deslocamento}
Para descrever mudanças geométricas em um meio contínuo é necessário fazer o mapeamento do deslocamento dos pontos dentro do domínio. O corpo em estudo é interpretado como um conjunto de partículas que inicialmente ocupam uma posição inicial $\mathbf{X}$, no momento $t_0$, que formam uma configuração de referência $C_0$ onde o corpo ainda não sofreu nenhum tipo de alteração. Após o corpo sofrer as alterações externas que mudaram sua geometria, as partículas ocupam novas posições em $\mathbf{x}$ no momento $t$ que formam uma configuração deformada $C$. O mapeamento do deslocamento é a equação que descreve a mudança de posição dessas partículas de um estado para o outro. Se a equação descreve a mudança do estado \textbf{indeformado} para o \textbf{deformado} é dado o nome de descrição \textit{Lagrangiana}, se ela descreve a mudança do estado \textbf{deformado} para o \textbf{indeformado} é dado o nome de descrição \textit{Euleriana}. 

O conjunto de pontos das partículas denominado por $\mathbf{X}$ na descrição \textit{Lagrangiana} é chamado de \textit{coordenadas materiais} e são usadas para descrever como o geometria de um corpo sólido muda após a aplicação de cargas externas, por esse motivo é usado na mecânica dos sólidos. O conjunto de pontos das partículas denominado por $\mathbf{x}$ na descrição \textit{Euleriana} é chamado de \textit{coordenadas espaciais} e são usadas para descrever como as propriedades dentro de um volume de controle mudam após a aplicação de condições de contorno, por esse motivo é usado na mecânica dos fluidos.

As equações que descrevem a mudança do grupo de coordenadas $\mathbf{X}$ para $\mathbf{x}$ pode ser representado por

\begin{equation}
	\mathbf{x} = \chi(\mathbf{X}, t)
\end{equation}

\begin{equation}
	\chi(\mathbf{X}, t_0) = \mathbf{X}
\end{equation}

%\begin{figure}
%	\centering
%	\includegraphics[width=0.6\linewidth]{fig_3-3-1_Reddy.png}
%	\caption{Mapeamento do Deslocamento}
%	\label{fig:my_label}
%\end{figure}

\section{Campo de Deslocamento}
O deslocamento dos pontos são a diferença entre as coordenadas das duas configurações do meio contínuo, como estamos falando de mecânica dos sólidos, daqui pra frente usaremos a descrição \textit{Lagrangiana} para montar as equações. Assim o deslocamento pode ser escrito como uma função das coordenadas \textit{materiais}

\begin{equation}
	\mathbf{u}(\mathbf{X}, t) = \mathbf{x}(\mathbf{X}, t) - \mathbf{X}
	\label{dispX}
\end{equation}

%\begin{figure}
%	\centering
%	\includegraphics[width=0.6\linewidth]{fig_3-3-3_Reddy.png}
%	\caption{Campo de Deslocamento}
%	\label{fig:my_label}
%\end{figu

\section{Tensor Gradiente de Deformação}
O tensor gradiente de deformação $\mathbf{F}$ é a relação entre a linha que liga as partículas $\mathcal{Q}$  e $\mathcal{P}$ separadas pela distância $d\mathbf{X}$ na configuração \textbf{indeformada}, e esta mesma linha agora com um tamanho $d\mathbf{x}$ na configuração \textbf{deformada}.  

\begin{equation}
	d\mathbf{x} = \mathbf{F} \cdot d\mathbf{X}
\end{equation}

\begin{equation}
	\mathbf{F} = (\frac{\partial\mathbf{x}}{\partial\mathbf{X}})^\mathbf{T} = (\nabla_{0}\mathbf{x})^T
	\label{Fdef}
\end{equation}

O operador gradiente $\nabla_{0}$ é escrito com o sub-índice 0 por se referir as derivadas de $d\mathbf{x}$ em relação a configuração indeformada $\mathbf{X}$, para escrever o gradiente em relação a configuração indeformada ele seria escrito como $\nabla_{\mathbf{x}}$. Desse modo o gradiente de deformação pode ser definido como o tensor transposto do gradiente das coordenas de uma partícula na configuração deformada $\mathbf{x}$ em relação as coordenadas na configuração indeformada $\mathbf{X}$.

\begin{equation}
	\mathbf{x} = \left[\begin{array}{ccc}
		x_1 \\ x_2 \\ x_3
	\end{array}\right],
	\qquad
	\mathbf{X} = \left[\begin{array}{ccc}
		X_1 \\ X_2 \\ X_3
	\end{array}\right]
\end{equation}

\begin{equation}
	\mathbf{F} = 
	\left[\begin{array}{ccc}
		\frac{\partial x_1}{\partial X_1} 
		& \frac{\partial x_1}{\partial X_2}
		& \frac{\partial x_1}{\partial X_3}
		\\
		\frac{\partial x_2}{\partial X_1} 
		& \frac{\partial x_2}{\partial X_2}
		& \frac{\partial x_2}{\partial X_3}
		\\
		\frac{\partial x_3}{\partial X_1} 
		& \frac{\partial x_3}{\partial X_2}
		& \frac{\partial x_3}{\partial X_3}
	\end{array}\right]
\end{equation}

Se estamos analisando duas partículas separadas por uma distância \textbf{finita} mudamos a notação $\partial$ por $d$.

Se reescrevermos a equação \ref{dispX} em termos de $\mathbf{x}$

\begin{equation}
	\mathbf{x} =  \mathbf{u} + \mathbf{X}
\end{equation}

aplicando o operador $\nabla_0$ em todos os termos temos

\begin{equation}
	\nabla_0 \mathbf{x} =  \nabla_0 \mathbf{u} + \nabla_0 \mathbf{X} 
	\label{Fdefnabla}
\end{equation}

lembrando a definição de $\mathbf{F}$ na equação \ref{Fdef} e aplicarmos a definição da equação \ref{Fdefnabla} podemos escrever um tensor em função dos deslocamentos e das coordenadas iniciais

\begin{equation}
	\mathbf{F} = (\nabla_0 \mathbf{u} + \nabla_0 \mathbf{X})^T
\end{equation}
  

Como $\nabla_0$ é uma derivada em relação as coordenadas iniciais, a operação $\nabla_0 \mathbf{X}$ resulta em uma matriz identidade

\begin{equation}
	\nabla_0 \mathbf{X} = \mathbf{I}
\end{equation}

\begin{equation}
	\nabla_0 \mathbf{X} = 
	\left[\begin{array}{ccc}
		\frac{\partial X_1}{\partial X_1}  = 1
		& \frac{\partial X_1}{\partial X_2} = 0
		& \frac{\partial X_1}{\partial X_3} = 0
		\\
		\frac{\partial X_2}{\partial X_1}  = 0
		& \frac{\partial X_2}{\partial X_2} = 1
		& \frac{\partial X_2}{\partial X_3} = 0
		\\
		\frac{\partial X_3}{\partial X_1} = 0
		& \frac{\partial X_3}{\partial X_2} = 0
		& \frac{\partial X_3}{\partial X_3} = 1
	\end{array}\right]
\end{equation}

assim podemos definir o tensor gradiente de deformação apenas em função dos deslocamentos

\begin{equation}
	\mathbf{F} = (\nabla_0 \mathbf{u} + \mathbf{I})^T
\end{equation}

%\begin{thebibliography}{}
%	\bibitem{reddy2009principles} Reddy, JN, (2009) Principles of continuum mechanics
%\end{thebibliography}
